%%% Hlavní soubor. Zde se definují základní parametry a odkazuje se na ostatní části. %%%

%% Verze pro jednostranný tisk:
% Okraje: levý 40mm, pravý 25mm, horní a dolní 25mm
% (ale pozor, LaTeX si sám přidává 1in)
\documentclass[12pt,a4paper]{report}
\setlength\textwidth{145mm}
\setlength\textheight{247mm}
\setlength\oddsidemargin{15mm}
\setlength\evensidemargin{15mm}
\setlength\topmargin{0mm}
\setlength\headsep{0mm}
\setlength\headheight{0mm}
% \openright zařídí, aby následující text začínal na pravé straně knihy
\let\openright=\clearpage

%% Pokud tiskneme oboustranně:
% \documentclass[12pt,a4paper,twoside,openright]{report}
% \setlength\textwidth{145mm}
% \setlength\textheight{247mm}
% \setlength\oddsidemargin{15mm}
% \setlength\evensidemargin{0mm}
% \setlength\topmargin{0mm}
% \setlength\headsep{0mm}
% \setlength\headheight{0mm}
% \let\openright=\cleardoublepage

\usepackage[british]{babel}
%% Pokud používáte csLaTeX (doporučeno):
%\usepackage{czech}
%% Pokud nikoliv:
%\usepackage[czech]{babel}
%\usepackage[T1]{fontenc}

%% Použité kódování znaků: obvykle latin2, cp1250 nebo utf8:
\usepackage[utf8]{inputenc}

%% Ostatní balíčky
\usepackage{graphicx}
\usepackage{amsthm}

%% Balíček hyperref, kterým jdou vyrábět klikací odkazy v PDF,
%% ale hlavně ho používáme k uložení metadat do PDF (včetně obsahu).
%% POZOR, nezapomeňte vyplnit jméno práce a autora.
\usepackage[ps2pdf,unicode]{hyperref}   % Musí být za všemi ostatními balíčky
\hypersetup{
	pdftitle={Properties of metric spaces by means of convergence},
	pdfauthor={Robin Pokorný}
%	, pdfkeywords={keyword1} {key2} {key3}
}

%%% Drobné úpravy stylu

% Tato makra přesvědčují mírně ošklivým trikem LaTeX, aby hlavičky kapitol
% sázel příčetněji a nevynechával nad nimi spoustu místa. Směle ignorujte.
\makeatletter
\def\@makechapterhead#1{
  {\parindent \z@ \raggedright \normalfont
   \Huge\bfseries \thechapter. #1
   \par\nobreak
   \vskip 20\p@
}}
\def\@makeschapterhead#1{
  {\parindent \z@ \raggedright \normalfont
   \Huge\bfseries #1
   \par\nobreak
   \vskip 20\p@
}}
\makeatother

% Toto makro definuje kapitolu, která není očíslovaná, ale je uvedena v obsahu.
\def\chapwithtoc#1{
\chapter*{#1}
\addcontentsline{toc}{chapter}{#1}
}

\begin{document}

% Trochu volnější nastavení dělení slov, než je default.
\lefthyphenmin=2
\righthyphenmin=2

% Definice a věty
\newtheorem{theorem}{Theorem}[chapter]
\newtheorem{define}[theorem]{Definition}	% Definice číslujeme
%\newtheorem*{define}{Definition} 	% Definice nečíslujeme


%%% Titulní strana práce

\pagestyle{empty}
\begin{center}

\large

Charles University in Prague

\medskip

Faculty of Mathematics and Physics

\vfill

{\bf\Large BACHELOR THESIS}

\vfill

\centerline{\mbox{\includegraphics[width=60mm]{logo.eps}}}

\vfill
\vspace{5mm}

{\LARGE Robin Pokorný}

\vspace{15mm}

% Název práce přesně podle zadání
{\LARGE\bfseries Properties of metric spaces by means of convergence}

\vfill

% Název katedry nebo ústavu, kde byla práce oficiálně zadána
% (dle Organizační struktury MFF UK)
Department of Mathematical Analysis

\vfill

\begin{tabular}{rl}

Supervisor of the bachelor thesis: & prof. RNDr. Miroslav Hušek, DrSc. \\
\noalign{\vspace{2mm}}
Study programme: & Mathematics \\
\noalign{\vspace{2mm}}
Specialization: & Obecná matematika \\
\end{tabular}

\vfill

% Zde doplňte rok
Prague 2012

\end{center}

\newpage

%%% Následuje vevázaný list -- kopie podepsaného "Zadání bakalářské práce".
%%% Toto zadání NENÍ součástí elektronické verze práce, nescanovat.

%%% Na tomto místě mohou být napsána případná poděkování (vedoucímu práce,
%%% konzultantovi, tomu, kdo zapůjčil software, literaturu apod.)

\openright

\noindent
Acknowledgement (?).

\newpage

%%% Strana s čestným prohlášením k bakalářské práci

\vglue 0pt plus 1fill

\noindent
I declare that I carried out this bachelor thesis independently, and only
with the cited sources, literature and other professional sources.

\medskip\noindent
I understand that my work relates to the rights and obligations under
the Act No. 121/2000 Coll., the Copyright Act, as amended, in particular the fact
that the Charles University in Prague has the right to conclude a license agreement
on the use of this work as a school work pursuant to Section 60 paragraph 1 of the Copyright Act.

\vspace{10mm}

\hbox{\hbox to 0.5\hsize{%
In ........ date ............
\hss}\hbox to 0.5\hsize{%
Signature
\hss}}

\vspace{20mm}
\newpage

%%% Povinná informační strana bakalářské práce

\vbox to 0.5\vsize{
\setlength\parindent{0mm}
\setlength\parskip{5mm}

Název práce:
Vlastnosti metrických prostorů pomocí konvergence
% přesně dle zadání

Autor:
Robin Pokorný

Katedra:  % Případně Ústav:
Katedra matematické analýzy
% dle Organizační struktury MFF UK

Vedoucí bakalářské práce:
prof. RNDr. Miroslav Hušek, DrSc., Katedra matematické analýzy
% dle Organizační struktury MFF UK, případně plný název pracoviště mimo MFF UK

Abstrakt:
% abstrakt v rozsahu 80-200 slov; nejedná se však o opis zadání bakalářské práce

Klíčová slova:
% 3 až 5 klíčových slov

\vss}\nobreak\vbox to 0.49\vsize{
\setlength\parindent{0mm}
\setlength\parskip{5mm}

Title:
Properties of metric spaces by means of convergence

Author:
Robin Pokorný

Department:
Department of Mathematical Analysis
% dle Organizační struktury MFF UK v angličtině

Supervisor:
prof. RNDr. Miroslav Hušek, DrSc., Department of Mathematical Analysis
% dle Organizační struktury MFF UK, případně plný název pracoviště
% mimo MFF UK v angličtině

Abstract:
% abstrakt v rozsahu 80-200 slov v angličtině; nejedná se však o překlad
% zadání bakalářské práce

Keywords:
% 3 až 5 klíčových slov v angličtině

\vss}

\newpage

%%% Strana s automaticky generovaným obsahem bakalářské práce. U matematických
%%% prací je přípustné, aby seznam tabulek a zkratek, existují-li, byl umístěn
%%% na začátku práce, místo na jejím konci.

\openright
\pagestyle{plain}
\setcounter{page}{1}
\tableofcontents

%%% Jednotlivé kapitoly práce jsou pro přehlednost uloženy v samostatných souborech
\chapter*{Preface}
\addcontentsline{toc}{chapter}{Preface}

In the theory of metric spaces various properties are studied. These properties usually have numerous characterisations. We shall examine whether selected ones can be described using convergent sequences.

In the first chapter we remind the definition of a metric space and some of its interesting properties without using sequences. Stated theorems show equivalent conditions or basic behaviour. 

The second chapter introduces two abstract spaces: sequential and uniformly sequential spaces. Both of them provide a generalization of certain properties of sequences we examined in the first chapter.  We show their relations to metric spaces and to each other.

In the third chapter properties are divided into three parts. Countinuous mappings, topology, compactness, connectedness and separability are those which can be characterised using sequential spaces only. We also examine separability behaves slightly differently althrought equivalently in metric spaces. Total boundedness and completeness need more than this, we show that can be characterised using uniformly sequential spaces. The last part concerns properties which cannot be characterised using either of these structures, we show boundedness is one of them.
\chapter{Title of the first chapter}

\section{Title of the first subchapter in the first chapter}

\section{Title of the second subchapter in the first chapter}


\chapter{Sequential and uniformly sequential spaces}

\section{Family of convergent sequences on a set}

To study the role of convergence in metric spaces we take some properties of sequences which hold in metric spaces and create new space using these properties as conditions. It turns out that the following conditions \emph{(C1)} - \emph{(C4)} introduce a rich class of spaces yet relatively simple to explore. 

\begin{define}\label{def:fcs}
Let $X$ be a non-empty set. A \emph{\fcs\ on $X$} (denoted by $\Xi_X$) is a set of pairs $(\{x_n\},x)$ consisting of a sequence $\{x_n\}\subseteq X$ and a point $x\in X$ satisfying the following conditions:
\begin{enumerate}
	\item[(C1)] If $(\{x_n\},x)\in\Xi_X$ and $\{y_n\}$ is a subsequence of $\{x_n\}$ then $(\{y_n\},x)\in\Xi_X$.
	\item[(C2)] If $(\{x_n\},x)\in\Xi_X$, $(\{x_n\},y)\in\Xi_X$ then $x=y$.
	\item[(C3)] If  $(\{x_n\},x)\notin\Xi_X$ then there exists a subsequence $\{x_{n_k}\}^\infty_{k=1}$ such that for none of its subsequence $\{x_{n_{k_i}}\}$ is $(\{x_{n_{k_i}}\},x)\in\Xi_X$.
	\item[(C4)] If $x\in X$ then $(\{x\},x)\in\Xi_X$.
\end{enumerate}
We say $(X,\Xi_X)$ is a \emph{sequential space}.
\end{define}

In the next we will write $\Xi$ instead of $\Xi_X$ wherever the set $X$ is clear from the context. Unless stated otherwise $X$ will denote a non-empty set. When $(\{x_n\},x)\in\Xi$ we say that $\{x_n\}$ is  a \emph{convergent sequence}, $\{x_n\}$ \emph{converges} to $x$ or that $x$ is a limit of  $\{x_n\}$.

The properties in Definition \ref{def:fcs} can be restated as follows:
\begin{enumerate}
	\item[(C1')] A subsequence of a convergent sequence is convergent and converges to the same point.
	\item[(C2')] Every sequence has at most one limit.
	\item[(C3')] A sequence that is not converging to a point $x$ contains a subsequence such that any of its subsequences either converges to a different point (not to $x$) or does not converge at all.
	\item[(C4')] A constant sequence converges to the common point.
\end{enumerate}
The first two and the fourth properties are natural and follow from basic knowledge of convergence in metric spaces. To understand the condition \emph{(C3)} we need to remind that a sequence $\{x_n\}$ which does not converge to $x$ might include a subsequence $\{x_{n_k}\}$ which converges to $x$.

Our motivation to define \fcss\ was to generalise convergence in metric spaces. The following shows the relation between them.

\begin{define}\label{def:gen}
Let $(X,\varrho)$ be a metric space. We say $\Xi$ is \emph{generated by} $X$ if 
\[
	\Xi=\{(\{x_n\},x): x\in X, \forall n\in \mathbb{N}, \ x_n\in X, \lim_{n \to \infty} \varrho(x_n,x)=0\}.
\]
\end{define}

A \fcs\ generated by a metric space is well defined. We know that $\lim_{n\to\infty}\varrho(x_n,x)=0$ implies $\lim_{n\to\infty}\varrho(x_{n_k},x)=0$ for every subsequence $\{x_{n_k}\}$ of $\{x_n\}$. If $x_n\to x$ and $x_n\to y$ for some $x,y\in X$ we have
\[
	0\le\varrho(x,y)\le\varrho(x_n,x)+\varrho(x_n,y)\to 0
\]
as $n\to\infty$, and so $\varrho(x,y)=0$ which by \emph{(M1)} means $x=y$. Property \emph{(C3)} follows from Theorem \ref{th:mpseq}. And last, obviously $\{x\}\to x$. 

In metric spaces adding a finite number of members to a sequence as well as removing finite number of members from it does not affect the convergence of the sequence. The same holds for \fcss.

\begin{theorem} \label{th:eqseq}
Let $\Xi$ be a \fcs, $(\{x_n\},x)\in\Xi$ and $a\in X$. Let $\{y_n\}$ and $\{z_n\}$ be sequences defined as follows:
\begin{enumerate}[(i)]
	\item $y_1=a$, $y_{n+1}=x_n\quad \forall n\in \mathbb{N}$,
	\item $z_n=x_{n+1}\quad \forall n\in \mathbb{N}$.
\end{enumerate}
Then $(\{y_n\},x), (\{z_n\},x)\in\Xi$.
\end{theorem} 
\begin{proof}
Let for contradiction suppose $(\{y_n\},x)\notin\Xi$. Then using \emph{(C3)} we can find a subsequence $\{y_{n_k}\}$ such that for every its subsequence $(\{y_{n_{k_j}}\},x)\notin\Xi$. We can assume $\{y_{n_{k_j}}\}$ does not include the first member of $\{y_n\}$. Then $\{y_{n_{k_j}}\}$ is a subsequence of $\{x_n\}$ and by \emph{(C1)} we have $(\{y_{n_{k_j}}\},x)\in\Xi$ which is a contradiction.

Since $\{z_n\}$ is a subsequence of $\{x_n\}$ the conclusion follows from \emph{(C1)}.
\end{proof}

\begin{theorem} \label{th:intxi}
An intersection of \fcss\ on a set $X$ is a \fcs\ on $X$.
\end{theorem} 
\begin{proof}
Let $\Xi_\alpha$ be a \fcs\ on $X$ for every $\alpha\in A$, where $A$ is an arbitrary non-empty set. Let $(\{x_n\},x)\in \bigcap_{\alpha\in A}\Xi_\alpha$ and $\{y_n\}$ be a subsequence of $\{x_n\}$. Then for every $\alpha\in A$ we have $(\{x_n\},x)\in \Xi_\alpha$ so by \emph{(C1)} we have $(\{y_n\},x)\in \Xi_\alpha$. That is, $(\{y_n\},x)\in \bigcap_{\alpha\in A}\Xi_\alpha$.

Also for every $\alpha\in A$ and every $x\in X$ we have $(\{x\},x)\in \Xi_\alpha$ so $(\{x\},x)\in \bigcap_{\alpha\in A}\Xi_\alpha$.

Let $(\{x_n\},x)\not\in \bigcap_{\alpha\in A}\Xi_\alpha$ then there exists $\alpha_0\in A$ for which $(\{x_n\},x)\not\in \Xi_{\alpha_0}$. We find the subsequence $\{x_{n_k}\}$ from \emph{(C3)} for $\Xi_{\alpha_0}$ and let $\{x_{n_{k_i}}\}$ be its subsequence. Then $(\{x_{n_{k_i}}\},x)\not\in\Xi_{\alpha_0}$ therefore $(\{x_{n_{k_i}}\},x)\not\in \bigcap_{\alpha\in A}\Xi_\alpha$.

Let $(\{x_n\},x)\in \bigcap_{\alpha\in A}\Xi_\alpha$ and $y\in X,\ x\ne y$. Then for some $\alpha\in A$ we have $(\{x_n\},x)\in \Xi_\alpha$ and from \emph{(C2)} for $\Xi_\alpha$ we have $(\{x_n\},y)\not\in \Xi_\alpha$. Therefore $(\{x_n\},y)\not\in \bigcap_{\alpha\in A}\Xi_\alpha$.
\end{proof}

\begin{define}\label{def:xi0}
Let $X$ be a non-empty set. Symbol $\Xi_0$ will denote the smallest (in the meaning of inclusion) \fcs\ on X.
%for which the following holds: \[\Xi_0\supset\bigg\{(\{x\},x) : x\in X\bigg\}.\]
\end{define}

%The existence of such \fcs\ follows from the Zorn's lemma and Theorem \ref{th:intxi}.
The definition says that $\Xi_0$ consists of convergent sequences which are constant up to a finite number of members.

Since $\Xi_0$ is the smallest \fcs\ on the set it is wherefore interesting to see for which metric spaces $\Xi=\Xi_0$ holds. It is easy to observe that $\Xi$ generated by a metrically discrete space (in Example \ref{ex:msp}) has his property.


At first sight one might think that  $\Xi=\Xi_0$ implies the space $(X,\varrho)$ to be metrically discrete. The following shows a counterexample constructing a space that is not metrically discrete.
\begin{example}
	Suppose $X=\{\frac{1}{n}, n\in \mathbb{N}\}$, then $X\subset \mathbb{R}$ and with induced metric it forms a subspace of $\mathbb{R}$. As the sequence $\{\frac{1}{n}\}$ has a limit in $\mathbb{R}$ which is not included in $X$, every convergent sequence in $X$ has to be constant from some index. Therefore $\Xi=\Xi_0$.
\end{example}

Another idea we had was: If $\Xi=\Xi_0$ then either $X$ is metrically discrete or $X$ is not complete. The following example prooves this statement to be false.

\begin{example}
	Suppose $X=\{(n,0), n\in \mathbb{N}\}\cup\{(n,\frac{1}{n}), n\in \mathbb{N}\}$, then $X\subset \mathbb{R}^2$. Intuitive image is the set of natural numbers as a subset of $x$ axis united with the graph of $\{\frac{1}{n}\}$ in the plane. This space (with induced metric) is complete and every convergent sequence is constant from some index.
\end{example}


\begin{theorem} \label{th:genxi0}
Let $\Xi$ be generated by a metric space $(X,\varrho)$ then the following are equivalent:
\begin{enumerate}[(i)]
	\item $\Xi=\Xi_0$;
	\item $\forall x\in X,\ \exists \varepsilon_x>0,\ \forall y\in X, x\neq y: \varrho(x,y)>\varepsilon_x$. (We say $(X,\varrho)$ is \emph{topologically discrete}.)
\end{enumerate}
\end{theorem} 

\begin{proof}
Let $X$ be a topologically discrete metric space, $\Xi$ generated by $X$ and for contradiction suppose there exists a sequence $\{x_n\}$ and $x\in X$ such that $(\{x_n\},x)\in\Xi\setminus\Xi_0$. Then from (ii) we find $\varepsilon_x$. Because $\lim_{n\to\infty}\varrho(x_n,x)=0$ there is $n_0\in\mathbb{N}$ for which $\varrho(x_{n_0},x)\le\frac{1}{2}\varepsilon_x$ and $x_{n_0}\ne x$ (since $(\{x_n\},x)\notin\Xi_0$). This is a contradiction with $X$ being topologically discrete.

Let $\Xi=\Xi_0$ and for contradiction suppose that for some $x\in X$ we have
\[
	\forall \varepsilon>0,\ \exists y\in X, x\ne y: \varrho(x,y)<\varepsilon.
\]
For $n\in \mathbb{N}$ we set $\varepsilon_n:=\frac{1}{n}$ and find $y_n\ne x$ such that $\varrho(x,y_n)<\varepsilon_n$. Then $\{y_n\}\to x$ which is a contradiction since $\{y_n\}$ is not constant.
\end{proof}

\section{Uniformly sequential structure on a set}

%\begin{define}{Properties needed later \tbd} \label{def:cofin}
%\begin{enumerate}[(i)]
%	\item $A'$ is cofinal in $A$.
%\end{enumerate}
%\end{define}

% \{x_n\} 

\begin{define} \label{def:unistr}
Let $X$ be a non-empty set.  A \emph{\uss\ on $X$} is an equivalence $\sim$ on sequences in $X$ satisfying the following conditions:
\begin{enumerate}
	\item[(U1)] If $\{x_n\} \sim \{y_n\}$ and $\{x_{n_k}\}, \{y_{n_k}\}$ are subsequences then $\{x_{n_k}\} \sim \{y_{n_k}\}$.
	\item[(U2)] If $x_n=x, y_n=y\ \forall n\in \mathbb{N}$ for some $x,y\in X$ and $\{x_n\} \sim \{y_n\}$ then $x=y$.
	\item[(U3)] If $\{x_n\} \not\sim \{y_n\}$ then there exist subsequences $\{x_{n_k}\}, \{y_{n_k}\}$ such that for none of their subsequences $\{x_{n_{k_i}}\}, \{y_{n_{k_i}}\}$ is $\{x_{n_{k_i}}\}\sim \{y_{n_{k_i}}\}$.
	\item[(U4)] If $x\in X$ then $\{x\} \sim \{y\}$. (Reflexivity) 
\end{enumerate}
We say $(X,\sim)$ is a \emph{uniformly sequential space}. If $\{x_n\} \sim\{y_n\}$ we say  \emph{$\{x_n\} $ is adjacent to $\{y_n\}$}.
\end{define}

\begin{theorem} \label{th:unitoseq}
	Let $X$ have a \uss. Then
\[
	\Xi_\sim:=\bigg\{(\{x_n\},x): \{x_n\}\sim \{x\}\bigg\}
\]
is a \fcs\ on $X$.
\end{theorem}
\begin{proof}
Let $\{x_n\}\sim \{x\}$ and $\{x_{n_k}\}$ be subsequence of $\{x_n\}$. We know every subsquence of a constant sequence is again a constant sequence. So from \emph{(U1)} we have $\{x_{n_k}\}\sim \{x\}$.

Let $x,y\in X$. Let $\{x_n\} \sim \{x\}$ and $\{x_n\} \sim \{y\}$. Then from transitivity of $\sim$ we obtain $\{x\} \sim \{y\}$ and \emph{(U2)} gives us $x=y$.

Let $\{x_n\} \not\sim \{x\}$. Then \emph{(C3)} follows directly from \emph{(U3)}.

Because of \emph{(U4)} we have $(\{x\},x)\in\Xi_\sim$.
\end{proof}

\begin{theorem}\label{th:adjprop}
Adjancent sequences in a metric space $(X,\varrho)$ satisfy conditions \emph{(U1)-(U4)}.
\end{theorem}
\begin{proof}
From Theorem \ref{th:adjquiv} we know  $\sim_\varrho$ is an equivalence.

Let $x_n=x, y_n=y\ \forall n\in \mathbb{N}$ for some $x,y\in X$ and $\{x_n\} \sim_\varrho \{y_n\}$. If $x\ne y$ then from \emph{(M1)} we have $\varrho(x,y)>0$ so $x=y$.

Let $\{x_n\} \sim_\varrho \{y_n\}$ and $\{x_{n_k}\}, \{y_{n_k}\}$ be subsequences. Then clearly $\{x_{n_k}\}\sim_\varrho \{y_{n_k}\}$ .

Let $\{x_n\} \not\sim_\varrho \{y_n\}$ then there exists $\varepsilon>0$ such that for infinitely many $n \in \mathbb{N}$ we have $\varrho(x_n,y_n)\ge\varepsilon$, these introduce a subsequence. It is easy to observe it has desired properties.
\end{proof}

We are not able to characterise a Cauchy sequence in terms of simple convergence. It turns out that this problem is the main motivation to develop uniformly sequential spaces in which we can define this property.

\begin{define}\label{def:ucauch}
Let $X$ have a \uss. We say a sequence $\{x_n\}$ in $X$ is \emph{Cauchy} if it is adjacent to all its subsequences.
\end{define}

As we would expect every convergent sequence is Cauchy.

\begin{theorem}\label{th:cauchconvg}
Let $X$ have a \uss. If $(\{x_n\}, x)\in\Xi_\sim$ then $\{x_n\}$ is Cauchy.
\end{theorem}
\begin{proof}
Let $\{x_{n_k}\}$ be a subsequence of $\{x_n\}$. Then from \emph{(C1)} we have $(\{x_{n_k}\}, x)\in\Xi_\sim$. That is $\{x_n\} \sim \{x\}$ and $\{x_{n_k}\} \sim \{x\}$ and the transitivity of $\sim$ gives us $\{x_{n_k}\} \sim \{x_n\}$.
\end{proof}

The relation of a Cauchy sequence in a uniformly sequential space and a Cauchy sequence in a metric space is shown later (Theorem \ref{th:scomplcompl}).

%Poznámka: (Více na Császar: GT, 1978 - str. 153 for totally bounded uniform spaces, str. 139 for uniform spaces)






















% Ukázka použití některých konstrukcí LateXu (odkomentujte, chcete-li)
%%% Ukázka použití některých konstrukcí LaTeXu

\subsection{Sample of \LaTeX}
\label{ssec:ukazka}

V~této krátké části ukážeme použití několika základních konstrukcí \LaTeX{}u,
které by se vám mohly při psaní práce hodit.

Třeba odrážky:

\begin{itemize}
\item Logo Matfyzu vidíme na obrázku~\ref{fig:mff}.
\item Tato subsekce má číslo~\ref{ssec:ukazka}.
\item Odkaz na literaturu~\cite{lamport94}.
\end{itemize}

Druhy pomlček:
červeno-černý (krátká),
strana 16--22 (střední),
$45-44$ (minus),
a~toto je --- jak se asi dalo čekat --- vložená věta ohraničená dlouhými pomlčkami.
(Všimněte si, že jsme za \verb|a| napsali vlnovku místo mezery: to aby se
tam nemohl rozdělit řádek.)

% Makro na české uvozovky (novější verze LaTeXu ho už mají zabudované)
%\uv{České uvozovky.}

\newtheorem{theorem}{Theorem}
\newtheorem*{define}{Definition}	% Definice nečíslujeme, proto "*"

\begin{define}
{\sl Strom} je souvislý graf bez kružnic.
\end{define}

\begin{theorem}
Tato věta neplatí.
\end{theorem}

\begin{proof}
Neplatné věty nemají důkaz.
\end{proof}

\begin{figure}
	\centering
	\includegraphics[width=30mm]{logo.eps}
	\caption{Logo MFF UK}
	\label{fig:mff}
\end{figure}


\chapter*{Conclusion}
\addcontentsline{toc}{chapter}{Conclusion}


%%% Seznam použité literatury
%%% Seznam použité literatury je zpracován podle platných standardů. Povinnou citační
%%% normou pro bakalářskou práci je ISO 690. Jména časopisů lze uvádět zkráceně, ale jen
%%% v kodifikované podobě. Všechny použité zdroje a prameny musí být řádně citovány.

\def\bibname{Seznam použité literatury}
\begin{thebibliography}{99}
\addcontentsline{toc}{chapter}{\bibname}

\bibitem{lamport94}
  {\sc Lamport,} Leslie.
  \emph{\LaTeX: A Document Preparation System}.
  2. vydání.
  Massachusetts: Addison Wesley, 1994.
  ISBN 0-201-52983-1.

\end{thebibliography}


%%% Tabulky v bakalářské práci, existují-li.
\chapwithtoc{List of Tables}

%%% Použité zkratky v bakalářské práci, existují-li, včetně jejich vysvětlení.
\chapwithtoc{List of Abbreviations}

%%% Přílohy k bakalářské práci, existují-li (různé dodatky jako výpisy programů,
%%% diagramy apod.). Každá příloha musí být alespoň jednou odkazována z vlastního
%%% textu práce. Přílohy se číslují.
\chapwithtoc{Attachments}

\openright
\end{document}
